\section{引言}\label{sec:introduction}
辐射源定位常先估计到达时间差等中间量，再由这些量求解位置。广义互相关与双曲线定位为这一流程提供了成熟工具\cite{KnappCarter1976,ChanHo1994}。当多个源使用重叠频带且波形未知时，各站接收的是不同传播响应下的混合信号，中间量的提取与跨站关联随之变难。直接位置估计（direct position determination，DPD）将传播关系直接带入观测模型，以位置作为主要待估参数，为多源联合处理提供了另一条途径\cite{ChenHudsonYao2002,WeissAmar2005}。

本文研究四个同步单天线站对两个未知宽带源的直接定位。两源占用相同频带，接收功率可以明显不同；定位器已知站点坐标及相对传播规律，不知道发射波形与发射功率。目标是在有限计算时间内准确恢复两个位置。单源空间谱中的强源及其旁瓣可能占据多个高分位置，直接取两个峰容易漏掉弱源；同时搜索两个二维位置则需要检查大量位置组合。逐源发现后再连续精修可以减少搜索量，但其效果取决于两个具体选择：保留哪些起点，以及怎样更新源位置。

已有直接定位方法从不同角度处理混合观测。宽带最大似然与交替投影（alternating projection，AP）固定部分位置、更新其余位置，降低联合搜索的维数\cite{ChenHudsonYao2002}。多重信号分类（MUSIC）型DPD根据候选响应与噪声子空间的正交性定位\cite{WeissAmar2005}；最小方差无失真响应（MVDR）利用协方差逆矩阵压低干扰\cite{Capon1969}，DPD-HR将这一思想用于直接定位\cite{TirerWeiss2016}。本文让这些方法使用相同观测、处理频带和传播信息，比较它们在强弱源混合时的实际表现。

稀疏参数估计将粗搜索与离网精修结合起来。Newtonized orthogonal matching pursuit（NOMP）在贪心发现后进行Newton更新和循环精修\cite{Mamandipoor2016}。多频率二维块NOMP已用于离网波束形成\cite{Yang2022}；QNOMP采用准Newton联合优化改善多径参数估计\cite{Zeng2025}。这些研究为多频率建模、二维精修和联合优化提供了基础。对于本文的多站直接定位，仍需根据位置相关的传播时延和幅度建立更新公式，并处理未知波形估计与位置搜索之间的耦合。本文重点研究：保留多个条件候选能否减少过早选择错误区域的情况，以及在局部位置更新中消去波形系数能否降低计算成本。

本文的求解思路是先保留可能的位置，再在同一拟合目标下充分比较。粗条件谱只评价网格点上的当前残差，排名第一的起点未必得到精修后的最佳解。因此，MCJ-AP保留多个条件候选，每个候选都与第一源位置一起优化，最后比较完整残差。相比反复进行全域搜索，这一流程只增加少量局部分支；相比逐源交替移动，联合更新允许两个位置同时调整。

未知波形的处理同样影响迭代效率。变量投影法先求最优线性系数，再优化非线性参数\cite{GolubPereyra1973}。这里线性系数是逐频波形，非线性参数是源位置。试探位置改变时重新估计波形，可以避免固定系数对位置步长的约束。本文据此构造波形消元块NOMP，并推导固定系数曲率与消元曲率的关系，解释局部更新为何可能需要不同的计算量。

围绕这些问题，本文形成以下三项工作。
\begin{enumerate}
  \item 针对同步单天线站的未知宽带强弱重叠源，建立一致的观测、功率和位置估计模型。通过波形消元，将条件搜索、残差发现和连续精修统一到位置最小二乘问题中。
  \item 提出多候选联合精修AP，将候选保留与联合位置更新结合起来。前者提高困难构型中的定位成功率，后者在相同候选与停止精度下降低计算成本；条件投影推导和完整配对实验分别说明其依据与作用。
  \item 构造波形消元块NOMP，用消元位置目标替代固定系数的局部更新。通过Schur补分析两种曲率及局部收敛差别，并以实际目标、梯度调用和耗时评价求解效率。
\end{enumerate}

下文依次给出观测模型、求解算法与计算量分析，再在20个源站布局、六档功率差下检验定位表现。全部2400条观测用于核心方法及消融，预先指定的240条共同观测用于八类方法比较。
